\documentclass[12pt,letterpaper]{article}
\usepackage[utf8]{inputenc}
\usepackage[T1]{fontenc}
\usepackage{times}
\usepackage[margin=0.5in,top=0.4in,bottom=0.4in]{geometry}
\usepackage{hyperref}
\usepackage{xcolor}
\usepackage{enumitem}
\usepackage{titlesec}

\hypersetup{colorlinks=true,linkcolor=blue!60!black,urlcolor=blue!60!black,citecolor=blue!60!black}

\setlength{\parindent}{0pt}
\setlength{\parskip}{2pt}
\setlist{nosep,leftmargin=*}

\titleformat{\section}{\normalsize\bfseries\scshape}{}{0pt}{}[\vspace{-2pt}\rule{\linewidth}{0.4pt}]
\titleformat{name=\section,numberless}{\large\bfseries\color{blue!60!black}}{}{0pt}{}[\vspace{-2pt}\rule{\linewidth}{0.6pt}]
\titlespacing{\section}{0pt}{8pt}{4pt}

\pagestyle{empty}

\begin{document}

{\footnotesize\textbf{Arxiv Digest --- 2026-11-05} \hfill 7 papers}

\smallskip

\section*{Multilingual NLP}

\begin{minipage}{\linewidth}
{\small\textbf{Familiar letters, not related languages, let unseen languages into a shared space of meaning}} {\scriptsize[\href{https://arxiv.org/abs/2611.06006}{arxiv}]}\\
{\scriptsize Yuliang Liu, Xinyu Wang, Xiang Bai}\\


{\footnotesize Many written languages lack translations, and it is unclear whether models can map an unseen language into a shared meaning space because it is related to a known language or because it uses familiar letters. The authors test this by pretraining GlyphMind, a tokenizer-free vision transformer that reads text as pixels, to align translations across 27 writing systems, then adapting it on untranslated Lao, whose language and script were unseen. GlyphMind retrieved English translations seven times as often as the strongest simple baseline. When Lao was rewritten in unfamiliar letters it failed to enter the shared space despite related Thai being in pretraining, while Santali entered in Devanagari but not in its own script.}

\end{minipage}

\medskip

\begin{minipage}{\linewidth}
{\small\textbf{Adapting Qwen3 to Tajik from 0.6B to 8B: Lower Text Loss Did Not Mean Better Understanding}} {\scriptsize[\href{https://arxiv.org/abs/2611.06062}{arxiv}]}\\
{\scriptsize Tohir Saidzoda, Madina Abdusalomzoda, Nasimjon Rajabov, Akhrorkhon Ikromov, Komron Giyosiev, Abubakr Abdullozoda, Amin Jalilov, Jonibek Kishvarov, Rahmonali Sharipov, Mariyamkhon Shodyeva}\\


{\footnotesize The authors test whether adapting Qwen3 to Tajik improves understanding, not just text loss. Small LoRA adapters trained on Tajik news and books reduced Tajik text loss across model sizes from 0.6B to 8B, but did not improve Tajik reading comprehension where measurable on Belebele. In contrast, at 8B, continued pretraining on about 375 million tokens of raw Tajik text followed by supervised fine-tuning produced a main model that scored 11.4 points above Qwen3-8B on Belebele, while English MMLU changed by -1.2 points. They also report documentation and release problems that limit interpretability and reproducibility.}

\end{minipage}

\medskip

\begin{minipage}{\linewidth}
{\small\textbf{ALTO: African and Levantine Tech-Facilitated Gender-Based Violence Corpus}} {\scriptsize[\href{https://arxiv.org/abs/2611.06706}{arxiv}]}\\
{\scriptsize Tianwei Zhang, Ivy Njoki Nyagwencha, Zahraa Dirani, Shalini Joshi, Scott A. Hale}\\


{\footnotesize Tech-facilitated gender-based violence is a pervasive online harm, but datasets and models for detecting gendered harm largely reflect Global North languages. Working with civil society organizations in Jordan, Francophone West Africa, and Kenya, the authors used a community-led process to crowdsource culturally and linguistically relevant examples in Levantine Arabic, African French, and Swahili, then applied active learning to annotate a dataset of 3,022 items called ALTO. Benchmarking and fine-tuning show that existing harmful-content models have poor recall for this harm on ALTO but improve after fine-tuning. Models often miss covert harms in all three languages, with markedly higher error rates in Swahili.}

\end{minipage}

\medskip

\begin{minipage}{\linewidth}
{\small\textbf{NaijaFacts: A Framework for Disentangling Language from Knowledge in Facts about Nigeria}} {\scriptsize[\href{https://arxiv.org/abs/2611.06504}{arxiv}]}\\
{\scriptsize Zainab Afolabi, Nikita Kozodoi, Jack Butler, Babs Khalidson}\\


{\footnotesize Multilingual benchmarks for African languages often translate generic questions, so they test language skill more than knowledge about Africa. The authors introduce NaijaFacts, a benchmark of factual and statistical questions about Nigeria written by native speakers and grounded in official statistics, with each question in English, Nigerian Pidgin, Yoruba, Igbo, and Hausa. By comparing similar questions for the UK and US, they separate a language effect from a knowledge effect, and find the knowledge effect dominates: English accuracy is 49--50\% for US/UK but 20\% for Nigeria. Even strong models reach 66--79\% on US/UK, yet none exceed 31\% on Nigerian statistics.}

\end{minipage}

\medskip


\section*{Multilingual Speech Processing}

\begin{minipage}{\linewidth}
{\small\textbf{Confucius4-R2T2: A Stable-Prefix Streaming ASR Model for Real Real-Time Transcription}} {\scriptsize[\href{https://arxiv.org/abs/2611.06043}{arxiv}]}\\
{\scriptsize Daniel Huang, Qi He, Shaoxiong Duan, Yingjie Li, Qiao Cheng, Yitao Duan}\\


{\footnotesize Streaming speech recognition systems often revise earlier words, which can disrupt translators or voice agents that act on the transcript immediately. Confucius4-R2T2 is trained to emit a stable opening portion of the transcript and leave it unchanged as more audio arrives. The authors obtain it by post-training Qwen3-ASR-1.7B on 1,500 hours of Chinese and English audio while freezing the audio encoder and projector. The method changes the training labels rather than the architecture. Evaluations show recognition quality close to the base model's full-audio output, with low emission delay and gains beyond the post-training languages.}

\end{minipage}

\medskip

\begin{minipage}{\linewidth}
{\small\textbf{Holistic Parallel Supervision for Expressive Speech to Speech Translation}} {\scriptsize[\href{https://arxiv.org/abs/2611.06010}{arxiv}]}\\
{\scriptsize Hao Zhang, Ismail Rasim Ulgen, Bismarck Bamfo Odoom, Rui Liu, Zihan Zhang, Berrak Sisman, Philipp Koehn}\\


{\footnotesize Expressiveness is central to speech-to-speech translation, but existing datasets rarely align source and target speech in content, prosody, timing, and speaker identity, so systems often depend on indirect conditioning. The authors propose a training-data construction strategy that pairs professional dubbing, which provides human-performed matches in content, prosody, and timing, with multilingual voice conversion to reduce speaker mismatch. They test this ``holistic parallel supervision'' on English--German and English--Spanish using existing expressive speech-to-speech translation systems and DirectS2ST, an end-to-end model that predicts the full target-speech representation from source speech. Holistic supervision improves expressive realization across systems, with DirectS2ST showing the largest gains, including higher human-rated expressiveness.}

\end{minipage}

\medskip


\section*{Foundation Model Theory}

\begin{minipage}{\linewidth}
{\small\textbf{Phonological Knowledge of Chinese Characters in Language Models}} {\scriptsize[\href{https://arxiv.org/abs/2611.05803}{arxiv}]}\\
{\scriptsize Xiaozhe Ji, Renfen Hu, Hai Hu}\\


{\footnotesize Language models' phonological knowledge of Chinese characters, and how they represent and use it, remains unclear. The authors introduce PinyinBench, a controlled benchmark covering homophones, phonological features (initials, finals, tones), and polyphonic characters, and evaluate 17 text-only and multimodal models using behavioral tests, layer-wise probing, and activation steering. A frontier model exceeds 90\% accuracy on homophone judgments when two characters share no phonological features, but drops to 72\% when they differ only in tone. Probing and steering suggest segmental information becomes more usable with depth, while tone is decodable but harder to influence generation; audio training improves decodability without improving steering.}

\end{minipage}

\medskip



\section*{Field Overview}
{\footnotesize Researchers are building and stress-testing language-model agents that can plan, use tools, and manage long histories, while developing faster inference methods and better benchmarks to measure memory, routing, and evaluation reliability. A second theme is safety, governance, and accountability: detecting misinformation, deepfakes, and machine-written text, auditing model updates and unlearning, and hardening agents against jailbreaks and policy failures. A third theme is broadening multimodal and multilingual capability, including speech and low-resource languages, and applying these models to domains like health, law, finance, and robotics.}


\end{document}